\scanpage{240}

IV -- \mspage{7} -- \hfill n°338

Avec les notations du th.2, l'ensemble $\{\underline p_0,\ldots,\underline p_{n-1}\}$ n'est pas nécessairement déterminé de façon unique par $M$ ; en particulier il peut être distinct de $\operatorname{Ass}(M)$.

PROPOSITION 8.- \uline{Soient} $A$ \uline{un anneau noethérien}, $\underline a$ \uline{un idéal de} $A$, \uline{et} $M$ \uline{un $A$-module de type fini}. \uline{Les conditions suivantes sont équivalentes} :

a) \uline{il existe un élément non nul} $x$ \uline{de} $M$ \uline{tel que} $\underline ax=(0)$,

b) \uline{pour tout} $a\in\underline a$, \uline{il existe un élément non nul} $x$ \uline{de} $M$ \uline{tel que} $ax=0$ ;

c) \uline{il existe} $\underline p\in\operatorname{Ass}(M)$ \uline{tel que} $\underline a\subset\underline p$.

Il est clair que a) implique b). D'après le cor. 2 de la prop.2, la condition b) signifie que $\underline a$ est contenu dans la réunion des idéaux premiers associés à $M$, donc dans l'un d'eux puisque $\operatorname{Ass}(M)$ est fini (Chap.II, §1, n°1, lemme 2) ; ainsi b) implique c). Enfin, s'il existe $\underline p\in\operatorname{Ass}(M)$ tel que $\underline a\subset\underline p$, $\underline p$ est l'annulateur d'un élément non nul $x$ de $A$ (déf.1), et on a $\underline ax=(0)$ ; donc c) implique a).

PROPOSITION 9.- \uline{Soient} $A$ \uline{un anneau noethérien}, $\underline a$ \uline{un idéal de} $A$, \uline{et} $M$ \uline{un $A$-module de type fini}. \uline{Pour qu'il existe un entier} $n>0$ \uline{tel que} $a^nM=(0)$, \uline{il faut et il suffit que $a$ soit contenu dans l'intersection des idéaux premiers associés à} $M$.

La condition est nécessaire en vertu du cor. de la prop.7. Réciproquement, ce même corollaire et le fait que $M$ est de type fini, montrent que tout élément $a$ de $\underline a$ définit une homothétie nilpotente de $M$ ; en particulier, si $(a_1,\ldots,a_n)$ est un système fini de générateurs de $\underline a$, il existe des entiers $n(i)>0$ tels que $a_i^{n(i)}M=(0)$ pour $1\leq i\leq q$ ; on a alors $\underline a^{n(1)+\cdots+n(q)}M=(0)$ ce qui termine la démonstration.

\begin{center}-:-:-:-\end{center}

\margin{239}
